\section{Introduction}
Steganography hides a message inside an innocent-looking cover so that only the intended receiver knows it is there.
In most deep steganography systems the cover is treated as disposable: the network adds a small perturbation, the
receiver reads the message, and the original image is gone for good. For medical scans, legal evidence, satellite and
military imagery this is not acceptable, because the exact original pixels may be needed later. Reversible data hiding
(RDH) solves the reversibility part. Classical RDH schemes based on difference expansion~\cite{tian2003de}, histogram
shifting~\cite{ni2006hs} and prediction-error expansion (PEE)~\cite{thodi2007pee,sachnev2009,ou2013ppee} return the cover
bit for bit, but the hidden bits are destroyed by the smallest change to the marked image, such as a single JPEG save.

Deep, and in particular invertible, networks have the opposite profile. Invertible steganography networks such as
HiNet~\cite{jing2021hinet} and ISN~\cite{lu2021isn} hide and reveal with the same weights, and noise-layer training as in
HiDDeN~\cite{zhu2018hidden} or RIIS~\cite{xu2022riis} makes the message survive distortion. However, these networks work in
floating point. When the stego image is rounded to 8 bits and the part of the output that is not transmitted (the
latent) is replaced by zeros, the inverse pass no longer gives back the cover. Recent robust-reversible work, DRRW and
CRMark~\cite{chen2025drrw,chen2025crmark}, closes this gap with an \emph{integer} invertible network: overflowed pixels and
the dropped latent are compressed into an auxiliary bitstream and hidden again with classical RDH.

We follow the same two-stage idea but change what has to be stored. Our observation is simple: if the coupling layers
are integer additions that are \emph{conditioned on the message}, then the receiver can undo them exactly as soon as it
knows the message, and there is no latent at all. The message itself is read back by a separate small decoder, and the
only things the receiver can get wrong are (i) a few message bits and (ii) pixels that were clipped at 0 or 255. These
two kinds of residual are what we compensate. Since the decoder outputs a confidence for every bit, the wrong bits are
concentrated where the confidence is low, and an arithmetic coder that uses this confidence as context codes them very
cheaply.

\textbf{Contributions.} The main contributions of this work are:
\begin{itemize}
\item a message-conditioned integer coupling network (IC-INN) whose inverse is exact by construction, so that
exact recovery needs no latent variable, only decoder-error flags and overflow corrections;
\item a residual compensation stage that codes these corrections with a context-adaptive binary arithmetic coder whose
contexts come from the decoder confidence and the local stego values, followed by a multi-layer rhombus PEE embedding
with an overflow-safe predictor range;
\item a lightweight design (62.9k parameters) that trains on a CPU, with a full MATLAB implementation of inference,
coding, embedding, extraction and an App Designer GUI;
\item an evaluation on 132 test images (BSD100, Kodak and real DICOM CT/MR/CR scans) against classical rhombus PEE and a
float-INN baseline with residual compensation, covering exact recovery, image quality, side-information size and
robustness to noise and JPEG.
\end{itemize}
